\documentclass[11pt]{amsart}

\usepackage[T1]{fontenc}
\usepackage{lmodern}
\usepackage{microtype}
\usepackage{mathtools,amssymb,amsthm}
\usepackage[hidelinks]{hyperref}

\hypersetup{
  pdftitle={Gap-Only Interlacing of the Real Roots of Consecutive Yablonskii-Vorob'ev Polynomials, Given the Published Real-Root Count}
}

\newtheorem{theorem}{Theorem}
\newtheorem{lemma}[theorem]{Lemma}
\theoremstyle{definition}
\newtheorem{definition}[theorem]{Definition}

\DeclareMathOperator{\sgn}{sign}
\newcommand{\Zs}{\mathcal Z}
\newcommand{\NN}{N}

\title[Gap-only interlacing of $Y_n$ and $Y_{n+1}$]
{Gap-Only Interlacing of the Real Roots of Consecutive Yablonskii--Vorob'ev
Polynomials, Given the Published Real-Root Count}

\date{}

\subjclass[2020]{34M55, 33E17, 26C10}
\keywords{Yablonskii--Vorob'ev polynomials, Painlev\'e II, real roots,
interlacing, Okamoto polynomials}

\begin{document}

\begin{abstract}
Let $Y_n$ be the Yablonskii--Vorob'ev polynomials. In the Conclusions of their
paper on the roots of generalised Okamoto polynomials, Roffelsen and Stokes
record that the numbers of real roots of the $Y_n$ are known and that
interlacing of the real roots of $Y_{n-1}$ and $Y_{n+1}$ is a theorem of
Clarkson, ``but interlacing of roots of consecutive polynomials $Y_n$ and
$Y_{n+1}$ has not been proved''. Their own definition of interlacing is a
two-sided condition on gaps alone: between two real roots of one polynomial
there lies a real root of the other. Granting three published facts about the
$Y_n$ --- simplicity of the roots, coprimality of consecutive members, and
Roffelsen's count of the real roots --- we prove the quoted statement for every
$n\ge1$ in that gap-only sense and in no other, and in the sharper form that the
merged real-root list of $Y_n$ and $Y_{n+1}$ has exactly $n+1$ elements,
strictly alternates, and begins with a root of $Y_{n+1}$. The qualification is
not idle. Under the stronger reading of interlacing that a reader may well have
in mind, which also demands that the two root counts differ by one, that same
merged word makes the
statement false for every odd $n$, since $m_{n+1}=m_n$ there; under that reading
nothing here settles the sentence quoted above, and we do not claim it does. A
sign identity read
off the Yablonskii--Vorob'ev recurrence forces the parity of the number of roots
of $Y_{n+1}$ in each gap between consecutive roots of $Y_n$, and the published
count of real roots leaves no room for more than the forced minimum.
\end{abstract}

\maketitle

\section{The statement}

The Yablonskii--Vorob'ev polynomials are defined by
\begin{equation}\label{eq:rec}
 Y_0=1,\qquad Y_1=z,\qquad
 Y_{n+1}Y_{n-1}=zY_n^2-4\bigl(Y_nY_n''-(Y_n')^2\bigr)\quad(n\ge1),
\end{equation}
and $Y_n\in\mathbb Z[z]$ is monic of degree $d_n=n(n+1)/2$. They carry the
rational solutions of the second Painlev\'e equation. Write
\[
 \Zs_k=\{x\in\mathbb R: Y_k(x)=0\},\qquad m_k=\#\Zs_k .
\]

The Conclusions of Roffelsen and Stokes \cite{RS}, quoted here from the e-print
\texttt{arXiv:2402.15887v1}, name one gap in the picture for the
Yablonskii--Vorob'ev family:

\begin{quote}
the numbers of real roots of these polynomials $Y_n$ are known
\cite{Roffelsen2012}, as well as interlacing of roots of $Y_{n-1}$ and
$Y_{n+1}$ on the real line \cite{ClarksonCMFT}, but interlacing of roots of
consecutive polynomials $Y_n$ and $Y_{n+1}$ has not been proved.
\end{quote}

\noindent
The two citations inside the quotation are \cite{Roffelsen2012} and
\cite{ClarksonCMFT} of the present bibliography. The journal version of record
was not consulted, so the wording above and the definition below are those of
that e-print:

\begin{definition}[\cite{RS}]\label{def:interlace}
The real roots of $P$ and $Q$ are \emph{interlaced} on a connected
$I\subseteq\mathbb R$ if between any two real roots of $P$ in $I$ there lies a
real root of $Q$, and between any two real roots of $Q$ in $I$ there lies a
real root of $P$.
\end{definition}

\noindent
This is a two-sided condition on gaps only; it does \emph{not} require the two
root counts to differ by one, and every assertion below is under it. Under a
stronger variant that adds that demand, the word proved below makes the
statement true for every even $n$ and false for every odd $n$, since
$m_{n+1}=m_n$ for odd $n$. So the sentence quoted above is settled here in the
gap-only reading only; in the stronger reading it is not settled but refuted for
odd $n$, and which of the two the quoted sentence intends we do not know.

\begin{theorem}\label{thm:main}
For every $n\ge1$ the real roots of $Y_n$ and $Y_{n+1}$ are interlaced on all
of $\mathbb R$ in the sense of Definition~\ref{def:interlace}. More precisely,
$\Zs_n\cup\Zs_{n+1}$ has exactly $n+1$ elements, its elements alternate
strictly between the two polynomials, and its least element belongs to
$\Zs_{n+1}$; for even $n$ its greatest element also belongs to $\Zs_{n+1}$.
Equivalently, writing $1$ for a root of $Y_{n+1}$ and $0$ for a root of $Y_n$
and reading left to right along the real line, the merged root word is
\[
 \underbrace{1010\cdots10}_{n+1}\quad(n\text{ odd}),
 \qquad
 \underbrace{1010\cdots01}_{n+1}\quad(n\text{ even}).
\]
\end{theorem}

\section{The inputs}

Everything below uses only \eqref{eq:rec} and the following published facts,
none of them ours:
\begin{itemize}
\item[(F1)] every root of $Y_n$ is simple;
\item[(F2)] $\gcd(Y_n,Y_{n+1})=1$;
\item[(F3)] $m_k=\lfloor(k+1)/2\rfloor$ for every $k\ge0$;
\item[(F4)] $Y_k$ is monic of degree $d_k=k(k+1)/2$, hence
 $\sgn Y_k(+\infty)=+1$ and $\sgn Y_k(-\infty)=(-1)^{d_k}$.
\end{itemize}
(F1) and (F2) are due to Fukutani, Okamoto and Umemura \cite{FOU}; see
\cite[Thm.~7.1]{ClarksonSMF} and \cite[Thm.~4]{Roffelsen2012}. (F3) was
conjectured by Clarkson \cite[Conj.~7.4]{ClarksonSMF} and proved by Roffelsen
\cite[Thm.~1]{Roffelsen2012}. (F4) is immediate from \eqref{eq:rec}.

\section{The proof}

\begin{lemma}\label{lem:sign}
For every $n\ge1$ and every $a\in\Zs_n$,
\[
 Y_{n+1}(a)\,Y_{n-1}(a)=4\,Y_n'(a)^2>0 .
\]
In particular $Y_{n+1}$ and $Y_{n-1}$ take the same, nonzero, sign at every real
root of $Y_n$.
\end{lemma}

\begin{proof}
Rearranging \eqref{eq:rec},
\begin{equation}\label{eq:l1}
 Y_{n+1}Y_{n-1}-4(Y_n')^2=Y_n\bigl(zY_n-4Y_n''\bigr),
\end{equation}
an identity in $\mathbb Z[z]$. The right-hand side vanishes at $a$, and
$Y_n'(a)\ne0$ because $a$ is a simple root by (F1).
\end{proof}

For an open interval $I$ whose endpoints (possibly $\pm\infty$) are not roots of
$Y_k$, write $\NN_k(I)=\#(\Zs_k\cap I)$. Since all roots of $Y_k$ are simple by
(F1) and non-real roots occur in conjugate pairs, $\NN_k(I)$ is odd if and only
if $Y_k$ takes opposite signs at the two ends of $I$.

\begin{lemma}[parity transfer]\label{lem:parity}
Let $n\ge1$ and $\Zs_n=\{a_1<\dots<a_{m_n}\}$ (nonempty, by (F3)). Put
$L=(-\infty,a_1)$, $G_i=(a_i,a_{i+1})$ for $1\le i<m_n$, and
$R=(a_{m_n},+\infty)$. Then
\begin{align*}
 \NN_{n+1}(G_i)&\equiv \NN_{n-1}(G_i)\pmod 2\quad(1\le i<m_n),\\
 \NN_{n+1}(R)&\equiv \NN_{n-1}(R)\pmod 2,\\
 \NN_{n+1}(L)&\equiv \NN_{n-1}(L)+1\pmod 2 .
\end{align*}
\end{lemma}

\begin{proof}
By (F2), applied to the pairs $(Y_{n-1},Y_n)$ and $(Y_n,Y_{n+1})$, no $a_i$ is a
root of $Y_{n-1}$ or of $Y_{n+1}$, so all four counts are over intervals with
non-root endpoints and the sign criterion above applies. By
Lemma~\ref{lem:sign}, $Y_{n+1}$ and $Y_{n-1}$ share their sign at each $a_i$;
by (F4) they are both positive at $+\infty$. Hence across each $G_i$, and
across $R$, the two polynomials change sign together or not at all, which is the
first two congruences. At $-\infty$, however, (F4) gives signs $(-1)^{d_{n+1}}$
and $(-1)^{d_{n-1}}$, and
\[
 d_{n+1}-d_{n-1}=\tfrac{(n+1)(n+2)}{2}-\tfrac{(n-1)n}{2}=2n+1
\]
is \emph{odd}, so those two signs are opposite while the signs at $a_1$ agree:
exactly one of $Y_{n+1},Y_{n-1}$ changes sign across $L$.
\end{proof}

\begin{proof}[Proof of Theorem~\ref{thm:main}]
For $k\ge0$ let $S(k)$ be the assertion
\begin{quote}
$\Zs_k\cup\Zs_{k+1}$ alternates strictly between $\Zs_k$ and $\Zs_{k+1}$, and
its least element lies in $\Zs_{k+1}$.
\end{quote}
By (F3), $\#(\Zs_k\cup\Zs_{k+1})=\lfloor\frac{k+1}2\rfloor+\lfloor\frac{k+2}2\rfloor=k+1$,
so $S(k)$ determines the merged word completely: it is $1010\cdots$ of length
$k+1$, which is the word displayed in Theorem~\ref{thm:main}. We prove $S(n)$
for all $n\ge0$ by induction.

$S(0)$ holds: $\Zs_0=\varnothing$ because $Y_0=1$, and $\Zs_1=\{0\}$.

Let $n\ge1$ and assume $S(n-1)$. Keep the notation of
Lemma~\ref{lem:parity}. Reading $S(n-1)$ off, the merged word of
$(\Zs_{n-1},\Zs_n)$ has length $n$, alternates, and begins with an element of
$\Zs_n$; therefore
\begin{equation}\label{eq:step1}
 \NN_{n-1}(L)=0,\qquad
 \NN_{n-1}(G_i)=1\ (1\le i<m_n),\qquad
 \NN_{n-1}(R)=[\,n\text{ even}\,],
\end{equation}
the last because the alternating word of length $n$ starting in $\Zs_n$ ends in
$\Zs_{n-1}$ exactly when $n$ is even. Lemma~\ref{lem:parity} now transfers these
to $Y_{n+1}$ as parities, hence as lower bounds:
\begin{equation}\label{eq:step2}
 \NN_{n+1}(L)\ \text{odd},\qquad
 \NN_{n+1}(G_i)\ \text{odd},\qquad
 \NN_{n+1}(R)\equiv[\,n\text{ even}\,]\pmod2 .
\end{equation}
Summing the forced minima over the $m_n+1$ pieces of $\mathbb R\setminus\Zs_n$,
\begin{equation}\label{eq:sum}
\begin{aligned}
 m_{n+1}
 &=\NN_{n+1}(L)+\sum_{i}\NN_{n+1}(G_i)+\NN_{n+1}(R)\\
 &\ge 1+(m_n-1)+[\,n\text{ even}\,]
 = m_n+[\,n\text{ even}\,].
\end{aligned}
\end{equation}
But (F3) makes the right-hand side \emph{equal} to $m_{n+1}$ in both parities,
because
$\lfloor\frac{n+2}2\rfloor=\lfloor\frac{n+1}2\rfloor+[\,n\text{ even}\,]$.
So \eqref{eq:sum} is an equality and every term is pinned at its minimum:
\[
 \NN_{n+1}(L)=1,\qquad \NN_{n+1}(G_i)=1,\qquad \NN_{n+1}(R)=[\,n\text{ even}\,].
\]
For odd $n$ the last value is forced to be $0$ rather than merely at most $1$
\emph{by its parity} in \eqref{eq:step2}. Thus
$\Zs_n\cup\Zs_{n+1}$ reads
\[
 \bullet<a_1<\bullet<a_2<\cdots<a_{m_n}\ (<\bullet\ \text{iff } n \text{ even}),
\]
with each $\bullet$ a single element of $\Zs_{n+1}$: it alternates strictly and
begins in $\Zs_{n+1}$. That is $S(n)$.

Finally, $S(n)$ gives Definition~\ref{def:interlace} on $I=\mathbb R$: between
any two real roots of $Y_n$ there lie two \emph{consecutive} ones, and
alternation puts a root of $Y_{n+1}$ between those; symmetrically with the roles
exchanged.
\end{proof}

\section{Related statements}

(F3) is load-bearing above: parity alone bounds each interval count from below,
and it is the published total that pins every one of them at its minimum. Our
use of the index-$2$ pair is confined to Lemma~\ref{lem:sign}, where only a sign
is needed.

For the index-$2$ pair $(Y_{n-1},Y_{n+1})$ an alternating chain is proved by
Roffelsen \cite[p.~6]{Roffelsen2012} and interlacing by Clarkson
\cite[Thm.~7.3]{ClarksonSMF}, \cite{ClarksonCMFT}; Roffelsen also gives
$\min\Zs_{n-1}>\min\Zs_{n+1}$, $\max\Zs_{n-1}<\max\Zs_{n+1}$
\cite[Thm.~2]{Roffelsen2012} and the separate counts of positive and negative
real roots \cite[Thm.~3]{Roffelsen2012}. Roffelsen and Stokes
\cite[Cor.~3.15]{RS} prove adjacent-direction interlacing for the generalised
Okamoto polynomials $Q_{m,n}$, $Q_{m,n-1}$.

\section{Verification}

A program \texttt{verify.py} accompanies this note as a control on the interior
of the induction, in exact integer and rational arithmetic and using only the
Python standard library. It generates $Y_0,\dots,Y_{43}$ from \eqref{eq:rec} by
exact division in $\mathbb Z[z]$ and checks the monicity and degrees (F4); the
identity \eqref{eq:l1}; (F1) and (F2) as squarefreeness and coprimality over
finite fields; the real-root count (F3) by exact isolation with a Fujiwara
bound for $n\le41$, where the expected value from (F3) is what certifies the
isolation complete, with an independent Sturm computation that uses (F3)
nowhere confirming the count only for $n\le12$; the occupancy statements
\eqref{eq:step1} and their forced counterparts for $Y_{n+1}$, for
$1\le n\le40$; and the arithmetic of \eqref{eq:sum} for $1\le n\le 200000$. Its
recorded run reports
\begin{center}
\texttt{VERDICT: ALL 24 CHECKS PASS}
\end{center}
Three of those checks, the ones whose names begin \texttt{regression-}, compare
the generated objects against coefficient lists, isolating intervals and merged
root words held as literals inside the program itself; those literals are not
displayed in this paper, and the checks are labelled in the run as the internal
regression checks they are. No check asserts agreement with a display, a table
or a numbered result that this paper does not contain. The program verifies no
published proof and it verifies the root order only for $n\le40$; the theorem is
the hand proof above and covers every $n$. It also confirms, for
$1\le n\le200000$, the arithmetic behind the failure of the stronger reading of
interlacing recorded after Definition~\ref{def:interlace}.

\begin{thebibliography}{9}

\bibitem{ClarksonSMF}
P.~A. Clarkson,
\emph{Special polynomials associated with rational and algebraic solutions of
the Painlev\'e equations},
in: Th\'eories asymptotiques et \'equations de Painlev\'e,
S\'emin. Congr. \textbf{14}, Soc. Math. France, Paris, 2006, pp.~21--52.
Section~7 contains Thm.~7.1 (simplicity), Thm.~7.3 (index-$2$ interlacing)
and Conj.~7.4 (the real-root count).

\bibitem{ClarksonCMFT}
P.~A. Clarkson,
\emph{Special polynomials associated with rational solutions of the Painlev\'e
equations and applications to soliton equations},
Comput. Methods Funct. Theory \textbf{6} (2006), no.~2, 329--401.
\href{https://doi.org/10.1007/BF03321618}{doi:10.1007/BF03321618}.

\bibitem{FOU}
S.~Fukutani, K.~Okamoto, and H.~Umemura,
\emph{Special polynomials and the Hirota bilinear relations of the second and
the fourth Painlev\'e equations},
Nagoya Math. J. \textbf{159} (2000), 179--200.
\href{https://doi.org/10.1017/S0027763000007405}{doi:10.1017/S0027763000007405}.

\bibitem{Roffelsen2012}
P.~Roffelsen,
\emph{On the number of real roots of the Yablonskii--Vorob'ev polynomials},
SIGMA Symmetry Integrability Geom. Methods Appl. \textbf{8} (2012), 099,
9~pp.
\href{https://doi.org/10.3842/SIGMA.2012.099}{doi:10.3842/SIGMA.2012.099};
\href{https://arxiv.org/abs/1208.2337}{arXiv:1208.2337}.

\bibitem{RS}
P.~Roffelsen and A.~Stokes,
\emph{On real and imaginary roots of generalised Okamoto polynomials},
J. London Math. Soc. \textbf{112} (2025), no.~4, e70329.
\href{https://doi.org/10.1112/jlms.70329}{doi:10.1112/jlms.70329};
\href{https://arxiv.org/abs/2402.15887}{arXiv:2402.15887}.
The sentence quoted in Section~1, Definition~\ref{def:interlace} and the
numbered results cited above are taken from \texttt{arXiv:2402.15887v1}
(24 February 2024).

\end{thebibliography}

\end{document}
